\section{Target multigrid와 source DF의 독립적 소유권}
\label{sec:target-multigrid}

\subsection{두 개의 선택 축}

Periodic DC-DFT 계산에는 서로 다른 역할을 하는 두 numerical 선택이 있다. 첫 번째는
source HF의 Coulomb/exchange와 orbital Hessian action을 계산하는
\textbf{source two-electron backend}다. FFTDF, GDF, RSDF 중 하나가 이 역할을 맡는다.
두 번째는 수렴된 HF density에서 target DFA의 $J$와 $E_{\mathrm{xc}}$를 평가하는
\textbf{target integrator}다. 두 선택을 하나의 ``DF method''로 합쳐 부르면 source와
target의 책임이 뒤섞인다.

이를 다음 ordered pair로 적는다.
\begin{equation}
  \mathcal R
  =\left(\mathcal B_{\mathrm{source}},\mathcal I_{\mathrm{target}}\right),
  \qquad
  \mathcal B_{\mathrm{source}}\in\{\mathrm{FFTDF},\mathrm{GDF},\mathrm{RSDF}\}.
  \label{eq:target-owner-split}
\end{equation}
Target integrator를 바꾸어도 source HF density $D^{\mathrm{HF}}$, HF energy, occupied/virtual
orbital과 source Hessian $A^{\mathrm{HF}}$는 변하지 않아야 한다. 바뀌는 것은
$E_{\mathrm{DFA}}[D^{\mathrm{HF}}]$의 target quadrature와 그로부터 만들어지는 target
Fock, orbital-gradient right-hand side 및 explicit target derivative다.

\subsection{세 target integration mode}

현재 GPU route는 다음 세 mode를 구분한다.

\begin{center}
\begin{longtable}{p{0.21\textwidth}p{0.22\textwidth}p{0.22\textwidth}p{0.23\textwidth}}
\toprule
Mode & Target $J$ & Target XC & Source HF/AX \\
\midrule
\texttt{uniform} & selected DF & uniform KNumInt & selected DF \\
\texttt{multigrid\_xc} & selected DF & multigrid & selected DF \\
\texttt{multigrid\_j\_xc} & multigrid & multigrid & selected DF \\
\bottomrule
\end{longtable}
\end{center}

\texttt{multigrid\_xc}를 M1이라 부르면 target effective potential은
\begin{equation}
  V_{\mathrm{eff}}^{\mathrm{M1}}[D]
  =J_{\mathcal B_{\mathrm{source}}}[D]
   +V_{\mathrm{xc}}^{\mathrm{MG}}[D]
  \label{eq:target-m1-veff}
\end{equation}
이다. 따라서 finite auxiliary basis를 사용하는 GDF/RSDF target $J$의 representation을
유지하면서 XC quadrature만 multigrid로 교체한다.

반면 \texttt{multigrid\_j\_xc}를 M2라 부르면
\begin{equation}
  V_{\mathrm{eff}}^{\mathrm{M2}}[D]
  =J_{\mathrm{MG}}[D]+V_{\mathrm{xc}}^{\mathrm{MG}}[D].
  \label{eq:target-m2-veff}
\end{equation}
M2의 $J_{\mathrm{MG}}$와 finite-auxiliary GDF/RSDF $J$는 서로 다른 수치
representation이다. Mesh와 auxiliary basis를 각각 수렴시키면 같은 Coulomb limit로
가지만, finite setting에서 energy가 같아야 할 이유는 없다. 따라서 M2는 단순한
GDF/RSDF drop-in acceleration이 아니라 별도 target model로 보고 비교해야 한다.

\subsection{Energy와 gradient ownership}

Mode에 따라 analytic gradient에서 target Coulomb derivative를 정확히 한 번만 더해야
한다. Ownership predicate를 $\chi_J(\mathcal I)$라 하면
\begin{equation}
  \mathbf g_{\mathrm{target}}
  =\begin{cases}
    \mathbf g_J^{\mathcal B_{\mathrm{source}}}
    +\mathbf g_{\mathrm{xc}}^{\mathrm{uniform}},
      &\mathcal I=\texttt{uniform},\\
    \mathbf g_J^{\mathcal B_{\mathrm{source}}}
    +\mathbf g_{\mathrm{xc}}^{\mathrm{MG}},
      &\mathcal I=\texttt{multigrid\_xc},\\
    \mathbf g_{J+\mathrm{xc}}^{\mathrm{MG}},
      &\mathcal I=\texttt{multigrid\_j\_xc}.
  \end{cases}
  \label{eq:target-gradient-owner}
\end{equation}
마지막 줄에서 selected-DF target $dJ/dR$를 다시 더하면 Coulomb response가 이중 계산된다.
반대로 M1에서 selected-DF $dJ/dR$를 생략하면 energy와 gradient가 서로 다른 Hamiltonian을
미분하게 된다.

Target hcore와 GTH pseudopotential은 세 mode 모두 selected one-electron provider가
소유한다. Multigrid object를 target KS object 전체에 암묵적으로 부착하면 upstream
\texttt{get\_hcore()} 또는 \texttt{get\_j()} dispatch까지 바뀔 수 있다. 따라서 구현은
low-level target XC/J 결과를 명시적으로 조합하고, hcore/PP ownership을 별도 metadata로
보존해야 한다.

\subsection{왜 target multigrid가 HF Hessian을 바꾸지 않는가}

DC-DFT의 adjoint equation은 source HF stationarity constraint를 미분한 것이다.
그러므로
\begin{equation}
  A_{\mathrm{HF}}^Tz=-g_{\mathrm{target}}
\end{equation}
에서 오른쪽 $g_{\mathrm{target}}$은 target integrator에 따라 달라지지만 왼쪽
$A_{\mathrm{HF}}$는 source backend를 유지한다. M1/M2를 선택했다고 해서 AX iteration의
$J/K$를 multigrid로 바꾸면, 더 이상 원래 HF density가 만족한 stationarity constraint의
response를 푸는 것이 아니다.

이 구분은 계산 흐름을 다음처럼 만든다.
\begin{enumerate}
  \item GDF 또는 RSDF로 source HF를 수렴시켜 $D^{\mathrm{HF}}$를 만든다.
  \item 선택한 target integrator로 $E_{\mathrm{DFA}}[D^{\mathrm{HF}}]$와
        $F_{\mathrm{target}}$을 만든다.
  \item $F_{\mathrm{target},vo}$에서 $g_{\mathrm{target}}$을 만든다.
  \item source GDF/RSDF AX로 $z$를 푼다.
  \item target explicit derivative와 source-HF response derivative를 조립한다.
\end{enumerate}

\subsection{Source mesh와 target mesh}

\texttt{cell.ke\_cutoff}로 resolve한 \texttt{cell.mesh}는 FFT 기반 source operator,
GTH local potential 및 여러 reciprocal-space 작업에 사용될 수 있다. Target multigrid
mesh는 target XC/J quadrature를 위한 별도 numerical parameter다. 두 mesh가 같은
기본값을 가질 수는 있지만 같은 정확도 의미를 갖지는 않는다.

Source와 target convergence를 분리해
\begin{equation}
  \lim_{\mathcal M_s\to\infty}
  \lim_{\mathcal M_t\to\infty}
  \mathbf g(\mathcal M_s,\mathcal M_t)
  \label{eq:target-grid-convergence}
\end{equation}
를 확인한다. 하나의 cutoff만 높이고 target mesh를 보고하지 않으면 어느
representation이 수렴한 것인지 알 수 없다. \texttt{target\_multigrid\_mesh=null}은
자동 memory cap을 뜻하지 않고 현재 cell mesh를 그대로 상속한다. 큰 cell에서는 resolved
triplet과 grid-point 수를 반드시 report해야 한다.

Multigrid는 grid를 여러 Gaussian/locality level로 나누지만 finite quadrature라는 사실은
변하지 않는다. 따라서 촘촘한 uniform grid보다 항상 정확하거나 translation invariant한
방법이라고 가정할 수 없다. Energy phase scan, translation sum과 full-gradient mesh
convergence를 각각 확인한다.

\subsection{Energy workspace 재사용과 수명}

Target energy에서 만든 multigrid task, sorted Gaussian pair와 target Fock은 gradient에서
재사용할 수 있다. Energy와 gradient가 서로 다른 multigrid object 또는 mesh를 사용하면
central energy의 derivative가 아니다. 반대로 target XC derivative가 끝난 뒤에도 큰 task
cache를 계속 보유하면 뒤따르는 GDF/RSDF response와 memory peak가 겹친다.

안전한 lifetime은 다음과 같다.
\begin{equation}
  \text{build at target energy}
  \longrightarrow
  \text{reuse at target gradient}
  \longrightarrow
  \text{release before DF response}.
  \label{eq:target-multigrid-memory}
\end{equation}
이는 grid point를 버리는 근사가 아니라, 마지막 consumer가 끝난 object의 소유권을
종료하는 memory 관리다.

\subsection{권장 사용 범위}

M1은 source HF와 finite-auxiliary target $J$를 유지한 채 target XC만 교체하므로 현재
가장 보수적인 opt-in route다. M2는 analytic하게 일관될 수 있지만 target $J$
representation도 바꾸므로 별도 convergence가 필요하다. PP-only multigrid는 central
hcore, explicit PP derivative, nuclear CPHF source 및 Z contraction을 모두 같은
representation으로 옮기기 전에는 사용할 수 없다. 일부 PP term만 바꾸면 energy와
gradient ownership이 불일치한다.
